\documentclass[11pt]{article}
\usepackage[margin=0.7in]{geometry}
\usepackage{titlesec}
\usepackage{enumitem}
\usepackage{parskip}

\titleformat{\section}{\Large\bfseries}{\thesection}{1em}{}
\titleformat{\subsection}{\large\bfseries}{\thesubsection}{1em}{}

\setlength{\parskip}{0.3em}

\title{\textbf{Week 7: Free Jazz (Early 1960s Onward)} \\ \large Quick Reference}
\author{}
\date{}

\begin{document}

\maketitle

\section*{Period Overview}

\textbf{What is Free Jazz?} Complete freedom from all traditional rules. No predetermined harmony (atonal), no steady pulse (free time), no fixed forms. Collective improvisation - all musicians improvising simultaneously with no hierarchy. Civil Rights era context: musical freedom = social/political freedom. Emotion over beauty, truth over pleasure. Key figures: Ornette Coleman, Cecil Taylor, Eric Dolphy, Albert Ayler, late Coltrane.

\textbf{Key Sound:} Chaotic, intense, rule-breaking, confrontational. No steady pulse, atonal, extended techniques (multiphonics, overblowing, screaming).

\textbf{Relation to Bebop:} COMPLETE REJECTION of bebop rules. Bebop had complex changes - free jazz has NO changes. Bebop had virtuosic lines - free jazz has extreme sounds. Bebop was art music - free jazz is CONFRONTATIONAL art. Emotional expression over technical mastery.

\section*{Required Listening (5 Pieces)}

\subsection*{Ornette Coleman - "Lonely Woman" (1959)}
This free jazz piece completely abandons predetermined harmony in favor of improvisation that follows emotional arc rather than chord changes. Coleman's plastic alto saxophone creates a buzzy raw timbre quality while Don Cherry plays pocket trumpet with muted nasal sound over a slow dirge-like tempo with loosely maintained pulse. The piece was recorded without piano to avoid any harmonic constraints, representing the opposite of bebop's complex rapid chord progressions. The haunting memorable melody and free harmonic approach exemplify Coleman's harmolodics theory where melody, harmony, and rhythm are equal. The album title "The Shape of Jazz to Come" announced this as a manifesto against bebop and hard bop conventions, connecting musical freedom to the Civil Rights era's social liberation movements.

\subsection*{Cecil Taylor - "Tales (8 Whisps)" from \emph{Unit Structures} (1966)}
This free jazz piece features percussive piano tone clusters where adjacent notes are played together, creating completely atonal textures that replace bebop's chord voicings. Taylor's New England Conservatory classical training is applied to destroy bebop conventions through gestural explosive playing at variable tempo with free time sections. Critics said he "doesn't swing" similar to early criticism of Monk, and the extremely loud dynamics represent confrontational art rather than bebop's art music approach. The classical composition influence is evident despite the chaotic surface, showing extended techniques applied to piano improvisation. Coltrane initially hated Taylor's music but by 1966 praised it, demonstrating the evolution of jazz beyond bebop toward complete harmonic and rhythmic freedom.

\subsection*{Eric Dolphy - "Straight Up and Down" from \emph{Out to Lunch} (1964)}
This piece features Dolphy on bass clarinet, an instrument he introduced to jazz, expanding instrumentation beyond bebop's standard alto and tenor saxophones. The medium tempo is more structured than pure free jazz, but the wide interval leaps and jagged unpredictable phrasing use extended techniques that exceed bebop's complexity. The vibraphone's shimmering quality adds to the angular melody construction and dissonant harmonic language. Dolphy extends bebop's virtuosity into new timbral and intervallic territory while maintaining some structural elements, making this transitional between modal and free jazz. DownBeat magazine called this "anti-jazz" in 1961, showing how free jazz's confrontational approach challenged conventional definitions of the music.

\subsection*{Albert Ayler - "Ghosts" from \emph{Spiritual Unity} (1964)}
This free jazz piece uses free time with no steady pulse, featuring Ayler's huge raw screaming tenor tone in extreme altissimo register that pushes the saxophone beyond bebop's technical limits. The trio format with no chordal instrument avoids any harmonic constraints, while a simple folk melody is exploded into extreme sounds through overblowing and multiphonics. Influenced by Coleman Hawkins' "Body and Soul" but rejecting bebop's complex melodic heads in favor of spiritual catharsis over intellectual complexity. The spiritual ecstatic energy connects free jazz's abandonment of musical rules to the Civil Rights era's demand for social freedom. This represents free jazz's emphasis on raw emotional truth and spiritual expression regardless of conventional beauty.

\subsection*{John Coltrane - "Expression" from \emph{Expression} (1967)}
This completely free piece abandons all predetermined harmony, form, and meter, representing Coltrane's final spiritual and musical statement before his death at age 40 in 1967. Coltrane's harsh extreme tenor tone uses multiphonics (playing multiple notes simultaneously) and overblowing screaming in dense chaotic textures with apocalyptic intensity. His journey from bebop's "Giant Steps" chord complexity through modal jazz's "A Love Supreme" to this free expression demonstrates jazz's progressive liberation from structural constraints. The complete abandonment of bebop's chord-based improvisation caused Coltrane to lose much of his audience with this radical turn toward confrontational art. This piece exemplifies free jazz's Civil Rights era context where complete musical freedom symbolized the demand for complete social and political freedom.

\section*{Key Terms}

\textbf{Atonality} - No tonal center \\
\textbf{Collective Improvisation} - All musicians improvising simultaneously, no hierarchy \\
\textbf{Harmolodics} - Coleman's theory: melody, harmony, rhythm equal \\
\textbf{Extended Techniques} - Multiphonics, overblowing, screaming \\
\textbf{Altissimo Register} - Extreme high notes beyond normal range \\
\textbf{Free Time} - No steady pulse or meter \\
\textbf{Plastic Saxophone} - Coleman's white plastic alto (couldn't afford brass) \\
\textbf{Tone Clusters} - Adjacent notes played together (Taylor) \\
\textbf{Multiphonics} - Multiple notes on one instrument \\
\textbf{Fire Music} - Archie Shepp's term for intense free jazz

\section*{Readings}

\textbf{Szwed, \emph{Jazz 101} Ch. 24} - Free Jazz \\
\textbf{Szwed, \emph{Jazz 101} Ch. 26} - The 1970s

\section*{Progressive Liberation (Weeks 5-7)}

\textbf{Week 5 HARD BOP:} Freed jazz from pure intellectualism by adding soul/groove. BUT KEPT: Complex harmony, virtuosity, forms. \\
\textbf{Week 6 MODAL JAZZ:} Freed jazz from constant chord changes by using modes. BUT KEPT: Structure (modes), pulse, beauty. \\
\textbf{Week 7 FREE JAZZ:} Freed jazz from ALL rules by abandoning harmony/form/meter. BUT KEPT: Improvisation, intensity, expression.

\section*{How Free Jazz Emerged from Modal Jazz}

Modal jazz still had structure (modes). Free jazz abandoned ALL predetermined structures including modes. Modal jazz maintained steady pulse - free jazz used free time. Modal jazz sought beauty/spirituality - free jazz sought raw emotional truth regardless of beauty. Civil Rights context: complete freedom in music = complete freedom in society.

\end{document}
